\begin{theorem}[riemannZetaResidue]
\label{riemannZetaResidue}
\lean{riemannZetaResidue}
% at /home/runner/work/trough/trough/library/PrimeNumberTheoremAnd/ZetaBounds.lean:174.0-192.6
\leanok
The Riemann zeta function $\zeta(s)$ has a simple pole at $s=1$ with residue $1$. In particular,
  the function $\zeta(s) - \frac{1}{s-1}$ is bounded in a neighborhood of $s=1$.
\end{theorem}
\begin{proof}
\uses{ResidueOfTendsTo}
\leanok
From \texttt{riemannZeta\_residue\_one} (in Mathlib), we know that
  $(s-1)\zeta(s)$ goes to $1$ as $s\to1$. Now apply Theorem \ref{ResidueOfTendsTo}.
  (This can also be done using $\zeta_0$ below, which is expressed as
  $1/(s-1)$ plus things that are holomorphic for $\Re(s)>0$...)
\end{proof}
